\documentclass[11pt]{article}
\usepackage[a4paper,margin=1in]{geometry}
\usepackage{amsmath,amssymb}
\usepackage{hyperref}
\title{PyTex Project Philosophy}
\author{PyTex Contributors}
\date{\today}

\begin{document}
\maketitle

\section{Purpose}

PyTex is intended to be a coherent Python environment for crystallographic texture and diffraction where scientific meaning is explicit, reproducibility is routine, and documentation quality is treated as part of the product.

\section{First Principles}

\begin{itemize}
  \item Frames, symmetry, and provenance are first-class data, not hidden assumptions.
  \item Stable public APIs should encode scientific meaning directly.
  \item Documentation should explain both \emph{what} a workflow does and \emph{why} the mathematics is valid.
  \item Teaching and research are peers, not separate products.
  \item Validation claims must be evidenced by tests, benchmarks, or traceable parity notes.
\end{itemize}

\section{Outcome}

The repository should make it difficult to write semantically ambiguous code and easy to construct rigorous, publication-quality workflows.

\end{document}
